%%%%%%%%%%%%%%%%%%%%%%%%%%%%%%%%%
%
%       EXERCÍCIO
%
%%%%%%%%%%%%%%%%%%%%%%%%%%%%%%%%%

\ifdefstring{\atividade}{prova}{%
    \ifdefstring{\modo}{objetivo}{%
        \renewcommand{\valorquestao}{\ValorQObj\ ponto}
    }{%
        \renewcommand{\valorquestao}{\ValorQDisc\ pontos}
    }%
}%

\begin{exercicioBanco}[\valorquestao]
A figura a seguir representa um quadrado com 24 cm de lado.

\begin{center}
    \begin{tikzpicture}[scale=0.4]

        % Pontos
        \coordinate (A) at (0,0);
        \coordinate (B) at (8,0);
        \coordinate (C) at (8,8);
        \coordinate (D) at (0,8);

        \coordinate (E) at (2,8);
        \coordinate (F) at (0,3);
        \coordinate (G) at (8,3);

        % Região hachurada
        \fill[pattern=north east lines]
        (A)--(B)--(G)--(E)--(F)--cycle;

        % Contornos
        \draw[thick]
        (A)--(B)--(G)--(E)--(F)--cycle;

        \draw[thick]
        (A)--(B)--(C)--(D)--cycle;

        % Letras
        \node[left]  at (A) {$A$};
        \node[right] at (B) {$B$};
        \node[above] at (C) {$C$};
        \node[above left] at (D) {$D$};

        % Medidas
        \draw[dashed] (D)--++(2,0);
        \draw[<->] (0,8.7)--(2,8.7);
        \node at (1,9.2) {$x$};

        \draw[dashed] (-0.6,8)--(-0.6,3);
        \draw[<->] (-0.6,3)--(-0.6,8);
        \node[left] at (-0.8,5.5) {$x+1$};

        \draw[dashed] (8.6,8)--(8.6,3);
        \draw[<->] (8.6,8)--(8.6,3);
        \node[right] at (8.8,5.5) {$5$};

    \end{tikzpicture}
\end{center}

Determine a expressão da área \(y\) da parte hachurada da figura.

% Define as alternativas
\newcommand{\alternativas}{%
    \begin{center}
        \begin{tabularx}{\textwidth}{XXX}
            (a) \(y=-x^{2}-x+456\). &
            (b) \(y=x^{2}+x+456\). &
            (c) \(y=-x^{2}+x+456\). \\[5pt]
            (d) \(y=-x^{2}-2x+456\). &
            (e) \(y=x^{2}-x+456\).
        \end{tabularx}
    \end{center}
}

% Define a resposta correta
\newcommand{\resposta}{A}

% Lógica condicional para exibição
\ifdefstring{\atividade}{lista}{%
    \alternativas
    \vspace{0.5em}
    
    \noindent\textbf{Resposta:} letra \textbf{\resposta}.
}{%
    \ifdefstring{\modo}{objetivo}{%
        \alternativas
    }{%
        % modo = discursiva → não mostra alternativas
    }
}
\end{exercicioBanco}

